%% Descente de charges : zone d'influence de dalle reprise par un ensemble poteaux-poutre.
%% Perspective cavalière (fuyante 30°, coefficient 0,5). La dalle (bleu) reçoit les charges
%% d'exploitation et de neige, les transmet à la poutre, aux trois poteaux, puis aux semelles
%% et au sol. Les sollicitations ne sont PAS nommées : c'est la question posée à l'élève.
\documentclass[tikz,border=4pt]{standalone}
\usepackage{liaisons}
\begin{document}
\begin{tikzpicture}[x=1cm,y=1cm,
  beton/.style={draw=solideB, line width=1pt, fill=solideB!25},
  charge/.style={-{Stealth[length=5pt]}, line width=0.9pt, solideA},
  cote/.style={line width=0.6pt, solideD, {Stealth[length=4pt]}-{Stealth[length=4pt]}},
  attache/.style={line width=0.4pt, black!45}]

\useasboundingbox (-0.9,-1.15) rectangle (8.2,4.75);

%% ---------------- dalle en perspective ------------------------------------
\filldraw[beton] (6.6,3.0) -- (6.6,3.4) -- (7.9,4.15) -- (7.9,3.75) -- cycle;  % joue droite
\filldraw[beton, fill=solideB!15] (0.6,3.4) -- (6.6,3.4) -- (7.9,4.15) -- (1.9,4.15) -- cycle; % dessus
\filldraw[beton] (0.6,3.0) rectangle (6.6,3.4);                                % face avant
\node[font=\small, text=solideB!70!black, anchor=west] at (0.75,3.18) {dalle};

%% ---------------- charges réparties sur la dalle ---------------------------
\foreach \x in {2.0,2.85,...,7.2} {\draw[charge] (\x,4.58) -- (\x,4.22);}
\node[font=\small, text=solideA, anchor=east, align=right] at (1.85,4.42)
     {charges $q_{ne}$\\\scriptsize exploitation + neige};

%% ---------------- zone d'influence et voiture ------------------------------
\node[font=\small, text=solideB!70!black, anchor=west] at (3.55,3.72)
     {zone d'influence : 348 m$^2$};
\begin{scope}[shift={(2.05,3.62)}, scale=0.85]
  \fill[black!18, draw=black!45, line width=0.5pt, line join=round]
    (0,0.10) -- (0.12,0.32) -- (0.62,0.32) -- (0.80,0.10) -- (1.05,0.08) -- (1.05,0) -- (0,0) -- cycle;
  \fill[black!45] (0.22,0.02) circle (0.075) (0.85,0.02) circle (0.075);
\end{scope}

%% ---------------- cote de l'épaisseur de dalle ------------------------------
\draw[attache] (0.6,3.0) -- (0.15,3.0) (0.6,3.4) -- (0.15,3.4);
\draw[cote] (0.28,3.0) -- (0.28,3.4);
\node[font=\scriptsize, text=solideD, anchor=east, inner sep=2pt] at (0.22,3.2) {0,40 m};

%% ---------------- poutre -----------------------------------------------------
\filldraw[beton, fill=solideB!45] (0.8,2.6) rectangle (6.4,2.85);
\node[font=\small, text=solideB!70!black, anchor=east] at (0.7,2.72) {poutre};
%% déformée (exagérée) de la poutre
\draw[line width=0.6pt, black!60, dash pattern=on 2pt off 1.8pt]
  plot[smooth, tension=0.8] coordinates {(1.2,2.6) (2.4,2.30) (3.6,2.6) (4.8,2.30) (6.0,2.6)};
\node[font=\scriptsize, text=black!60, anchor=north] at (2.4,2.26) {déformée (exagérée)};

%% ---------------- poteaux ----------------------------------------------------
\foreach \x in {1.2,3.6,6.0} {\filldraw[beton] (\x-0.175,0.3) rectangle (\x+0.175,2.6);}
\node[font=\small, text=solideB!70!black, anchor=west] at (1.45,1.85) {poteau};
%% efforts dans un poteau (sollicitation à identifier)
\draw[charge] (3.6,1.95) -- (3.6,1.55);
\draw[charge] (3.6,0.75) -- (3.6,1.15);
\node[font=\scriptsize, text=black!60, anchor=west] at (3.85,1.35) {sollicitation ?};

%% ---------------- semelles et sol ---------------------------------------------
\foreach \x in {1.2,3.6,6.0} {\filldraw[draw=solideE, line width=1pt, fill=solideE!25]
  (\x-0.5,0) rectangle (\x+0.5,0.3);}
\node[font=\scriptsize, text=solideE!60!black, anchor=west] at (1.82,0.14) {semelle};
\draw[line width=0.9pt] (-0.3,0) -- (7.6,0);
\foreach \s in {-0.25,-0.05,...,7.5} {\draw[line width=0.45pt, black!70] (\s,0) -- (\s-0.16,-0.20);}
\draw[cote] (3.1,0.15) -- (4.1,0.15);
\node[font=\scriptsize, text=solideD, anchor=west, inner sep=2pt] at (4.15,0.15) {0,80 m};
%% réactions du sol
\foreach \x in {1.2,3.6,6.0} {\draw[charge] (\x,-0.72) -- (\x,-0.06);}
\node[font=\scriptsize, text=solideA, anchor=west] at (6.25,-0.68) {réaction du sol};

%% ---------------- cote de la hauteur des poteaux --------------------------------
\draw[attache] (6.175,0.3) -- (7.0,0.3) (6.175,2.6) -- (7.0,2.6);
\draw[cote] (6.85,0.3) -- (6.85,2.6);
\node[font=\scriptsize, text=solideD, anchor=west, inner sep=2pt] at (6.9,1.45) {4,75 m};
\end{tikzpicture}
\end{document}
